\documentclass[10pt,a4paper]{amsart}
\usepackage[margin=19mm]{geometry}
\usepackage{amsmath,amssymb,mathtools}
\usepackage[scr=rsfs,cal=euler]{mathalfa}
\usepackage{xfrac}
\usepackage[unicode,hidelinks,bookmarks=false]{hyperref}
\newcommand{\ie}{i.e.\ }
\newcommand{\cf}{cf.\ }
\newcommand{\resp}{respectively\ }
\newcommand{\Subsection}{\subsection}
\providecommand{\nobreakdash}{\nobreak\hbox{-}}
\setlength{\emergencystretch}{2em}
\multlinegap0pt
\usepackage{amssymb}
\usepackage[shortlabels]{enumitem}
\usepackage{cancel}
\usepackage{tikz}
\newtheorem{lemma}{Lemma}
\newcommand{\C}{\mathbb{C}}
\newcommand{\R}{\mathbb{R}}
\renewcommand{\leq}{\leqslant}
\title{Convex functions with symplectic Hessian}
\author{Jose Rafael Santiago Arellano}
\date{Source: arXiv:2608.23236v1 (CC BY 4.0)}
\begin{document}
\maketitle

\noindent\emph{Source and licence.} Convex functions with symplectic Hessian. Version arXiv:2608.23236v1 (CC BY 4.0), distributed by arXiv under the Creative Commons Attribution 4.0 International licence, http://creativecommons.org/licenses/by/4.0/, as recorded in the arXiv licence metadata for this exact version and shown on the abstract page https://arxiv.org/abs/2608.23236v1. Full licence notice: \texttt{licenses/arxiv-2608.23236-CC-BY-4.0.txt}.\par\smallskip

\noindent\textit{Editorial scope.} Counted unit: the article's lemma 3NN (source0.tex lines 499--503). The article's complete original proof of it is delivered with it (source0.tex lines 505--552, one \texttt{proof} environment of 48 lines). No mathematics is written, repaired or re-derived: the statement and the proof are the article's own lines, in the article's own notation, and no retained line is substituted or re-flowed.  Retained uncounted as the source-local context the proof needs: lines 337--339; lines 381--383; lines 424--426; lines 433--435; lines 466--472.  Nothing was omitted from any retained span, so the package carries no omission record.  No retained line contains a citation, so the wrapper carries no bibliography and no bibliography entry.  No retained line contains a figure environment, an image-inclusion command or drawing code, so no image is delivered and the package declares no asset dependency.  Editorial formatting only, in addition: an AMS \texttt{amsart} preamble that restates the article's own declarations and macros used by the retained lines, namely the environments \texttt{lemma} on separate counters, together with the standard packages those lines need.  The retained environments print in this document as Lemma 1 in order of appearance, whereas the article prints the same items with the numbering of its own full document.  The complete native source of the article as distributed in the arXiv e-print for this exact version is bundled unchanged as \texttt{sources/208400/source0.tex}.
\begin{equation}\label{symplectic}
u(O)=0, \quad (\nabla u)(O)=0, \quad  \operatorname{Hess}(u)|_{O}=I_{2m}, \quad \operatorname{Hess}(u)\in{\rm Sp}(2m,\R).
\end{equation}

\begin{lemma}\label{Zholo}
    In our setting, $Z := \Gamma+\sqrt{-1}\Phi$ is holomorphic in the coordinates $\alpha_i$.
\end{lemma}

\begin{equation}\label{T3m}
     T_{3}(x_{i},y_{i}) := T_{1}(x_{i},y_{i})+T_{2}(x_{i},y_{i})= (u_{x_{i}},y_{i})+(x_{i},u_{y_{i}})= (x_{i}+u_{x_{i}},y_{i}+u_{y_{i}}).
\end{equation}

\begin{align}\label{def_sigma}
\sigma_{i} \circ T_3:=(x_{i}+u_{x_{i}})+\sqrt{-1}(y_{i}+u_{y_{i}}).
 \end{align}

\begin{lemma}\label{Holchart1}
    The vector-valued map
    \begin{align}\label{def_f}f:T_{3}(B_{\delta(O)}(O))\to \C^{m}, \quad f_{i} \circ T_3 :=(x_{i}-u_{x_{i}})-\sqrt{-1}(y_{i}-u_{y_{i}}),\end{align}
    is holomorphic with respect to the coordinates $\sigma_i$. Moreover,
    $$\frac{\partial f}{\partial \sigma}=\mathscr{C} \circ Z \circ T_1 \circ T_3^{-1},$$
    where $Z: T_1(B_{\delta(O)}(O)) \to \mathfrak{H}_m$ is as in Lemma \ref{Zholo} and $\mathscr{C}: \mathfrak{H}_m \to \mathbb{D}_m$ is the Cayley transformation.
\end{lemma}

\begin{lemma}\label{3NN}
Let $u:U\subset\R^{2m}\to\R$ be smooth and convex with symplectic Hessian such that \eqref{symplectic} holds. As before, let $\delta(O)$ be the maximal radius such that $B_{\delta(O)}(O)\subset U$. Consider the map $$T_{3}:B_{\delta(O)}(O)\to T_{3}(B_{\delta(O)}(O))$$
from \eqref{T3m}. Let $R>0$ be the maximal radius such that $B_R(O)\subset T_{3}(B_{\delta(O)}(O))$. Then
$$\frac{R}{4}\leq \delta(O)\leq \frac{3R}{4}.$$
\end{lemma}

\begin{proof}
By Lemma \ref{Holchart1} and \eqref{symplectic} and by the definition of $\mathscr{C}$, we have that 
\begin{equation}\label{op}
    f(O) = 0, \quad \frac{\partial f}{\partial \sigma}(O) = \mathscr{C}(\sqrt{-1}I_m) = 0, \quad I_{m}-\frac{\partial f}{\partial \sigma}\,\overline{\frac{\partial f}{\partial \sigma}}^{\top}>0.
\end{equation}

\noindent {\bf Claim:} Let $\|\cdot\|_{\operatorname{op}}$ denote the operator norm on $\C^{m\times m}$ defined by 
$$\|A\|_{\operatorname{op}}:=\sup_{|v|=1} |Av|,$$ 
where $|\cdot|$ is the standard Euclidean norm on $\C^{m}$. Then we have that 
$$\left\| \frac{\partial f}{\partial \sigma}(\sigma)\right\|_{\operatorname{op}}\leq \frac{|\sigma|}{R}\;\,\text{for all}\;\,\sigma \in B_R(O).$$

\noindent {\bf Proof of the Claim:}
The Claim is trivial from \eqref{op} at $\sigma = O$. Thus, fix any $\sigma\in B_{R}(O)\setminus\{O\} $. For any two unit vectors $v,w\in \C^{m}$, define an auxiliary function $\phi_{v,w}:B_{R}(0)\subset\C\to \C$ by
$$\phi_{v,w}(\zeta):= \left\langle \frac{\partial f}{\partial \sigma}\left(\zeta \frac{\sigma}{|\sigma|}\right) \cdot v,w \right\rangle,$$
where $\langle \cdot, \cdot \rangle$ is the standard Hermitian inner product on $\mathbb{C}^m$ (complex linear in the first variable).
This is holomorphic because $f$ is. Moreover, $\phi_{v,w}(0) = 0$ and $|\phi_{v,w}(\zeta)| < 1$ for all $\zeta \in B_R(0)$ because \eqref{op} implies that all of the singular values of $\partial f/\partial \sigma$ are $< 1$, so that $\|\partial f/\partial \sigma\|_{\rm op} < 1$. Thus, the Schwarz lemma yields $|\phi_{v,w}(\zeta)| \leq |\zeta|/R$. The Claim follows by choosing $\zeta = |\sigma|$ and maximizing over $v,w$. \hfill $\Box$\medskip

\noindent To proceed, note that since $f(O)=0$,
$$
  |f(\sigma)| = \displaystyle \left|\int_{0}^{1} \frac{\partial f}{\partial \sigma}(t\sigma)\sigma \, dt\right|  
           \leq  \displaystyle \int_{0}^{1} \left\|\frac{\partial f}{\partial \sigma}(t\sigma)\right\|_{\operatorname{op}}|\sigma| \, dt 
           \leq  \displaystyle  \int_{0}^{1} t\frac{|\sigma|}{R}|\sigma| \, dt 
           =  \displaystyle  \frac{|\sigma|^2}{2R} 
           <   \displaystyle \frac{R}{2}
$$
if $|\sigma|<R$. Therefore, by using the equation
\begin{align*}
    T_3^{-1}(\sigma) = x+\sqrt{-1}y=\frac{\sigma+\overline{f(\sigma)}}{2},
\end{align*}
which is obvious from \eqref{def_sigma}, \eqref{def_f}, we conclude that
\begin{equation}\label{xy-sigma}
\frac{|\sigma|}{2}-\frac{R}{4} < |T_3^{-1}(\sigma)| < \frac{|\sigma|}{2}+\frac{R}{4} \;\,\text{for all}\;\,\sigma \in B_R(O).
\end{equation}

\noindent {\bf Proof of the inequality $R/4\leq\delta(O)$:} Let $0<\varepsilon<R$, and consider $B_{R-\varepsilon}(O)\subset T_{3}(B_{\delta(O)}(O))$. Set $$U_{\varepsilon}:=T_{3}^{-1}(B_{R-\varepsilon}(O))\subset B_{\delta(O)}(O).$$
Then $U_\varepsilon$ contains the point $O = T_3^{-1}(O)$, and $U_\varepsilon$ is an open set because $T_3$ is continuous. Moreover, $\overline{U_{\varepsilon}}\subset B_{\delta(O)}(O)$ as well because for all $u \in \partial U_\varepsilon$ we can pick a sequence $u_i \in U_\varepsilon$ converging to $u$; then $v_i := T_3(u_i) \in B_{R-\varepsilon}(O)$, so the sequence $v_i$ has a limit $v \in \overline{B_{R-\varepsilon}(O)} \subset B_R(O) \subset T_3(B_{\delta(O)}(O))$ after passing to a subsequence, so $u = T_3^{-1}(v) \in B_{\delta(O)}(O)$ by the continuity of $T_3^{-1}$.

Let $\delta_{\varepsilon}(O)$ denote the Euclidean distance from $O$ to $\partial U_{\varepsilon}$. Since $\overline{U_\varepsilon}\subset B_{\delta(O)}(O)$, there exists $z_\varepsilon\in\partial U_\varepsilon $ such that $|z_{\varepsilon}|=\delta_{\varepsilon}(O) < \delta(O)$. If  $\sigma_{\varepsilon}:=T_{3}(z_{\varepsilon})$, then $\sigma_\varepsilon\in \overline{B_{R-\varepsilon}(O)}$ as above. In fact, $\sigma_\varepsilon \in \partial B_{R-\varepsilon}(O)$ because otherwise $z_\varepsilon \in T_3^{-1}(B_{R-\varepsilon}(O)) = U_\varepsilon$. Thus, by \eqref{xy-sigma} we have that
$$\frac{R-\varepsilon}{2} - \frac{R}{4} = \frac{|\sigma_{\varepsilon}|}{2}-\frac{R}{4} < |z_{\varepsilon}| = \delta_\varepsilon(O) < \delta(O).$$
Letting $\varepsilon\to 0$, we conclude the claim.\medskip

\noindent {\bf Proof of the inequality $\delta(O)\leq 3R/4$:} Let $0<\varepsilon<\delta(O)$, and consider $B_{\delta(O)-\varepsilon}(O)\subset B_{\delta(O)}(O)$. Set $$V_{\varepsilon}:=T_{3}(B_{\delta(O)-\varepsilon}(O))\subset T_{3}(B_{\delta(O)}(O)).$$
Then $V_\varepsilon$ contains the point $O = T_3(O)$, and $V_\varepsilon$ is an open set because $T_3^{-1}$ is continuous. Moreover, $\overline{V_{\varepsilon}}\subset T_3(B_{\delta(O)}(O))$ as well because for all $v \in \partial V_\varepsilon$ we can pick a sequence $v_i \in V_\varepsilon$ converging to $v$; then $u_i := T_3^{-1}(v_i) \in B_{\delta(O)-\varepsilon}(O)$, so the sequence $u_i$ has a limit $u \in \overline{B_{\delta(O)-\varepsilon}(O)} \subset B_{\delta(O)}(O)$ after passing to a subsequence, so $v = T_3(u) \in T_3(B_{\delta(O)}(O))$ by the continuity of $T_3$.

Let $R_{\varepsilon}$ denote the Euclidean distance from $O$ to $\partial V_{\varepsilon}$. Since $\overline{V_{\varepsilon}}\subset T_{3}(B_{\delta(O)}(O))$, there exists $\sigma_{\varepsilon}\in \partial V_{\varepsilon}$ such that $|\sigma_{\varepsilon}|=R_{\varepsilon} < R$. If $z_{\varepsilon}:=T_{3}^{-1}(\sigma_{\varepsilon})$, then $z_\varepsilon\in \overline{B_{\delta(O)-\varepsilon}(O)}$ as above. In fact, $z_\varepsilon \in \partial B_{\delta(O)-\varepsilon}(O)$ because otherwise $\sigma_\varepsilon \in T_3(B_{R-\varepsilon}(O)) = V_\varepsilon$. Thus, by \eqref{xy-sigma} we have that
$$\delta(O)-\varepsilon = |z_{\varepsilon}|< \frac{|\sigma_{\varepsilon}|}{2}+\frac{R}{4} = \frac{R_\varepsilon}{2}+\frac{R}{4} < \frac{3R}{4}.$$
Letting $\varepsilon\to 0$, we conclude the claim.
\end{proof}

\end{document}
